\section{Problema de investigación (máximo 1/2 página)}
	
En los últimos años, las pequeñas y medianas empresas como minimarkets y bodegas enfrentan un severo problema de robo en tiendas, que causa pérdidas de inventario y merma económica en el sector retail peruano. Gran parte de la seguridad con la que cuentan dichos establecimientos depende únicamente de sistemas de videovigilancia tradicionales (CCTV), los cuales presentan varias limitaciones operativas; entre ellas, están la fatiga visual de los operadores humanos, la presencia de puntos ciegos y su naturaleza aleatoria. Al mismo tiempo, han surgido diversas técnicas de aprendizaje profundo (Deep Learning) como una alternativa novedosa y prometedora; sin embargo, la literatura científica revela tres vacíos que es importante resaltar:

\begin{enumerate}[leftmargin=*, noitemsep, topsep=2pt] 
		\item \textbf{Falta de modelado de conductas pre-delictivas:} La mayoría de los modelos se centran unicamente en detectar el momento exacto en donde ocurre el robo de forma reactiva, ignorando lo más importante, que seria los patrones de movimiento, los cuales caracterizan y son de gran ayuda para identificar la intención delictiva de la persona.
\item \textbf{Elevado costo computacional:} Las arquitecturas actuales requieren de hardware costoso, como servidores y GPUs de alto rendimiento, ocasionando que los comercios minoristas locales no puedan adquirir dicho hardware debido a su presupuesto limitado.
\item \textbf{Escasez de investigación en el contexto nacional:} Existe una gran escasez de estudios y benchmarks adaptados a la infraestructura, iluminación y ángulos de cámara del retail peruano, identificándose unicamente el antecedente de \cite{aguero_intelligent_2025}.
	\end{enumerate}
